What comprises the bank business?

the acceptance of funds from others as deposits or of other
unconditionally repayable funds from the public, unless the
claim to repayment is securitised in the form of bearer or
order bonds, irrespective of whether or not interest is paid
(deposit business)

the granting of money loans and acceptance credits (credit
business)

the purchase of financial instruments at the credit
institution's own risk for placing in the market or the
assumption of equivalent guarantees (underwriting
business)”

They should be differentiated from financial service institutions
(e.g. investment banks, venture capitalist funds) and 
financial undertakings (special purpose vehicles).


What are the transformation functions of a Bank?

Lot Size transformation
▪ Take many smaller deposits, lend larger amounts

Maturity transformation
▪ Deposits have very short maturities (e.g. can be withdrawn any time). 
However, given the many depositors, most of the time there‘s enough capital 
for long term loans.
▪ Banks may issue long-term bonds and grant short-term overdraft loans to customers.

Major Risks:
▪ Liquidity risk
▪ Interest rate risks

Risk transformation
Rather impressive: Banks turn risky loans into almost
riskless deposit liabilities.

A bank does so by:
▪ Diversification
▪ Using Equity as buffer


Derive the profit of banks in perfect markets
and the optimization problems of investors/firms.

The setup:
INSERT SLIDE 22

Assumptions are
Perfect Financial Market:
No frictions such as taxes or transaction costs
Goods are infinitely divisible and tradable
You can lend and borrow any amount of money
informationally efficient
Agents maximize \( \mathbb{E}[u(X)] \)

The investors problem is
\( \text{max} u(C_1, C_2) \)
s.t. 
\( C_1 + B_l + D^{+} = W_1 \)
\( C_2 = \pi_F + \pi_B (1 + r)B_l + (1 + r_D)D^{+} \)
Profit by Firm \( \pi_F \) and Bank \( \pi_B \) 
Securities earn interest at \( r \).
therefore deposits must also earn at least the same amount \( r_D = r \).

The firms problem
\( \text{max} \pi_F = f(I) - (1+r)B_f - (1+r_L)L^{-} \)
s.t. \( I = B_F + L^{-} \)
Firm maximizes profit \( \pi_F \) which depends on production function \( f \).
Issued Securities must be paid back with interest rate \( r \)
Loans must be paid back with interest rate \( r_L \)
Since Loans and Securities are riskless \( r = r_L \).

The banks problem

\( \text{max}\pi_B = r_L L^{+} - rB_B - r_D D^{-} \)
s.t. \( L^{+} = B_B + D^{-} \)
Bank maximizes profit \( \pi_B \)
Issued Securities must be paid back with interest rate \( r \)
Deposits must be paid back with interest rate \( r_D \), \( \implies \) \( r = r_D \)


The equilibrium solution in the perfect market is
\( I = S \)
\( D^{+} = D^{-} \)
\( L^{+} = L^{-} \)
\( B_l = B_F + B_B \)
Therefore Banks make zero profit! More importantly, investor‘s
or firm‘s decision are not affected by banks.


State the axioms of expected utility theory
and \( CE, RP \).

Def. of expected utility
In the discrete case
\( \mathbb{E}[u(X)] = \sum_{i=1}^n \mathbb{P}(\omega_i)u(X(\omega_i)) \)
otherwise
\( \mathbb{E}[u(X)] = \int_\Omega u(X)d\mu \)

Complete Ordering
• Completeness:
\( \forall X, Y, X \geq Y \) or \( Y \geq X \).
• Transitivity:
\( X \geq Y \text{and} Y \geq Z \implies X \geq Z \).
Continuity
if \( X \geq Y \geq Z \), then there exists some probability \( p \)
such that: \( b \sim pX + (1-p)Z \)
Independence
if \( X \geq Y \) it must hold for each lottery \( Z \) and probability \( p \).
\( pX + (1-p)Z \geq pY + (1-p)Z \).

Certainty equivalent 
Let it be denoted by \( c = CE(X) \) then it satisfies
\( u(c) = \mathbb{E}[u(X)] \).

Risk premium
Let the initial wealth level be \( w \) then the risk premia \( RP(X, w) \) must satisfy 
\( u(w + \mathbb{E}[X] - RP) = \mathbb{E}[u(w + X)] \).

Tabularize the Arrow Pratt risk attitudes and corresponding utility functions
and define ARA, RRA.

\begin{tabular}{|l|l|l|l|}
\hline
Arrow Pratt risk attitude & Utility function & \( RP = EV - CE \) & Risk attitude \\ \hline
\( ARA(x) = 0, \forall x > 0 \) & Linear: \(u_{\text{RN}}\) & \( RP(X) = 0, \forall X \) & Risk neutral \\ \hline
\(ARA(x) > 0 \forall x > 0\) & Concave: \(u_{\text{RA}}\) & \(RP(X) > 0 \forall X\) & Risk averse \\ \hline
\(ARA(x) < 0 \forall x > 0\) & Convex: \(u_{\text{RP}}\) & \(RP(X) < 0 \forall X\) & Risk prone \\ \hline
\end{tabular}

\begin{tabular}{|l|l|l|l|}
\hline
Arrow Pratt risk attitude & Utility function & $ RP = EV - CE $ & Risk attitude \\ \hline
$ ARA(x) = 0, \forall x > 0 $ & Linear: $u_{\text{RN}}$ & $ RP(X) = 0, \forall X $ & Risk neutral \\ \hline
$ARA(x) > 0 \forall x > 0$ & Concave: $u_{\text{RA}}$ & $RP(X) > 0 \forall X$ & Risk averse \\ \hline
$ARA(x) < 0 \forall x > 0$ & Convex: $u_{\text{RP}}$ & $RP(X) < 0 \forall X$ & Risk prone \\ \hline
\end{tabular}

The Arrow/Pratt measure of risk aversion
\( \text{ARA}(x) = - \frac{u''(x)}{u'(x)} \)
\( \text{RRA}(x) = x\text{ARA}(x)\)


What are the 2 important classes of utility functions?


CARA (constant absolute risk aversion)
\( \text{ARA}(x) = \gamma, \forall x \)
Which is equivalent to
\( u(x) = \alpha + \beta e^{-\gamma x}, \beta < 0 \)
It can be interpreted as
CARA investors do not become more risk averse or risk prone if their levels of wealth
change (or equivalently: if we add some fix amount of money to every lottery outcome).
For an initial level of wealth \( w \) and amount of money \( a \), lottery \( l \) it holds
\( \text{max}_a \mathbb{E}[u(w - a + al)] = u(W) \text{max}_a \mathbb{E}[u(-a + al)] \),
so the optimal choice of investment does not depend on the wealth.

CRRA (constant relative risk aversion)
\( \text{RRA}(x) = c, \forall x \)
\( u(x) = \alpha + \beta g(x) \) where
\[
g(x) = \begin{cases}
\frac{x^{1-\gamma}}{1 - \gamma} & \text{if} \gamma \neq 1 \\
\ln(x) & \text{else}
\end{cases} 
\]
For a relative amount \( aw \) of investment the optimum satisfies

\[
\text{max}_a \mathbb{E}[u(w((1-a) + al))] = \begin{cases}
u(w)\text{max}_{a}\mathbb{E}[u((1-a) + al)] & \text{if} \gamma \neq 1 \\
u(w) + \text{max}_{a}\mathbb{E}[u((1-a) + al)]& \text{else}
\end{cases} 
\]
Therefore the invested amount \( aw \) is proportional to the investor‘s wealth \( w \).

How is the repayment function for an optimal contract with symmetric information
determined?

Risky earnings \( \tilde{y} \) from project.
Observable to both parties (lender and borrower)

Optimal contract between lender and borrower is determined if we specify repayment function \(R(\tilde{y})\)
Repayment \(R(\tilde{y})\) to lender given some realization \(\tilde{y}\)
Lender requires minimum utility level \(U_L^0\)
Borrower retains \(\tilde{y} - R(\tilde{y})\)

Thus, solve:
\[
\begin{aligned}
    & \max_{R(\cdot)} Eu_B(\tilde{y} - R(\tilde{y})) \\
    & \text{s.t.} \\
    & \quad Eu_L(R(\tilde{y})) \geq U_L^0 \\
    & \quad 0 \leq R(\tilde{y}) \leq \tilde{y}
\end{aligned}
\]
with the help of the Lagrangian we derive
\( 𝐿 = 𝑢_𝐵(𝑦 − 𝑅) + 𝜆(𝑢_𝐿(𝑅) − 𝑈_𝐿^0 ) \)
\( R'(y) = \frac{\text{ARA}_B(y - R(y))}{\text{ARA}_B(y - R(y)) + \text{ARA}_L(R(y))} \)

For a given probability space, this differential equation may admit a solution,
the initial conditions are then determined via 
\( \mathbb{E}[u_L(R(\tilde{y}))] = \sum_i \mathbb{P}[\omega_i]u(R(\omega_i) = U^{0}_L \).


Give examples of asymmetric information

Ex ante uncertainty
Hidden information
Hidden characteristic

Ex interim uncertainty
Moral Hazard (firm exerts less effort than promised)
Hidden action 

Ex post uncertainty
Costly state verification


Give an example for agency costs

Agency costs are costs that arise when there are conflicts of
interest between stakeholders.
Since managers often hold shares of the firm's equity, a conflict
of interest exists if investment decisions have different
consequences for the value of equity and the value of debt

% TODO do this with a negative starting value
% such that the initial debt is partially repaid for the safe option
% and compare it to a risky bet with outcomes (completely paid vs. (more debt)), 
% to handle the leveraged case.

As a concrete example, suppose the firm's equity is \( 30$ \)
and there are the following return options
\( \mathbb{P}[\omega_1] = \mathbb{P}[\omega_2] = 0.5\)
\[
\forall \omega \in \Omega, X_1(\omega) = 60\$
\]
\[
X_2 = \begin{cases}
    0\$ & \omega_1 \\
    100\$ & \omega_2
\end{cases}
\]
Supposedly the upfront cost for both is \( \$40 \), now suppose the company takes on dept
to finance this endeavor at an interest rate of \( 10\% \) (this doesn't really matter for the example).
Suppose 
Then for equity holders the expected payoffs are
\( \pi_E(X_i) = \mathbb{E}[\max(X_i - 44, 0)] \)
thus \( \pi_E(X_1) = 16, \pi_E(X_2) = 28 \).
Now for the lenders this payoff is
\( \pi_L(X_i) = \mathbb{E}[\text{min}(44\%, 30\% + X_i)] - 40 \)
thus \( \pi_L(X_1) = 4, \pi_L(X_2) = -3 \).

When a firm faces financial distress, shareholders can gain by
making sufficiently risky investments, even if they have
negative NPV.



As a concrete example, suppose the firm's own equity is \( 30\$ \)
and debt \( -50\$ \) payable to the next period.
The firm can choose one of two investments, with the following profile
\( \mathbb{P}[\omega_1] = \mathbb{P}[\omega_2] = 0.5\)
\[
\forall \omega \in \Omega, X_1(\omega) = 50\$
\]
\[
X_2 = \begin{cases}
    0\$ & \omega_1 \\
    100\$ & \omega_2
\end{cases}
\]
Supposedly the upfront cost for both is \( 30\$ \).
Then for equity holders the expected payoffs are
\( \pi_E(X_i) = \mathbb{E}[\max(X_i - 50, 0)] \)
thus \( \pi_E(X_1) = 0, \pi_E(X_2) = 50\).

Now for the lenders the payoffs are
\( \pi_L(X_i) = \mathbb{E}[\text{min}(50\$, X_i)] - 50 \)
thus \( \pi_L(X_1) = 0, \pi_L(X_2) = -25\).

When a firm faces financial distress, shareholder can gain by making
sufficiently risky investments, even if they have NPV.


Summarize the debt overhang problems (too high leverage ratio)

Underinvestment problems by
▪ Foregoing positive NPV project (agency cost example)
▪ Disinvesting
▪ Excessive payout
    Stock buy-back programs instead of dividends
    wages for manager-owned firms
    Selling assets and immediately turning them into dividens (and not payback loans)
Overinvestment problems by
▪ Asset substitution (riskier projects)
▪ Executing negative NPV projects


How can financial strength be signaled with debt
Assume
\( \mathbb{P}[\omega_1] = \mathbb{P}[\omega_2] = 0.5\)
and let the firms value be \( V(\omega_1) = 100, V(\omega_2) = 50 \).
Suppose the equity owners know due to insider information that \( \omega_1 \)
will happen with absolute certainty.

Of course, debt holders know about these problems
Thus, they demand discounts on debt claims, e.g. they pay less
for bonds
Debt holders try to discipline managers by preferring short term
over long term debt
More likely that equity holders gain advantage over debt
holder in case of long term debt rather than short term debt
Also: multi-stage financing with debt

TODO those points are not really clear.


When should debt be forgiven, restructured or let be?

Restructuring in this context is meant as extending maturity (in general restructuring
also refers to debt forgiveness).

Company is in distressed because of negative demand shocks
→ restructure
Defaulting on debt due to debt overhang and resulting
underinvestment or excessive payout is a strategic default
→ beneficial to forgive debt to solve debt overhang problems

However the kind of signals to use, to differentiate between liquidity default and strategic default
are not obvious.
Low prior cash holdings and prior operating performance are not great indicators.
(Snow example with Ski hotels in Australia)


Why are SEOs a bad signal

If managers/insiders can assess true equity value better than
outside investors, they are likely to sell equity when equity is
overvalued.

SEO (Seasoned equity offerings) is an implicit signal that stock prices are too high.
Thus, stock price typically decreases when SEOs are
announced.
To reduce the adverse selection / lemmon problem, firms
prefer to announce SEO when there is least asymmetric
information, i.e. after earnings announcements, for example

Implications for banks in Financial Crisis
▪ High uncertainty in bank assets
▪ Raising additional equity (SEO) is bad signal
▪ Combination of high uncertainty and large SEO scares off existing
and new shareholders because
▪ equity likely to be overvalued in management‘s view, i.e.
more/higher future losses expected
▪ future losses to be shared with new shareholders
▪ Consequence: devaluation of existing equity→ more severe crisis!
▪ Solution:
▪ forced equity infusion by regulator/central bank or
▪ open books (create symmetric information) to sovereign
wealth funds or large private investor

Elaborate on the free banking era in the US (1837-1864)

Anyone could open a bank: no charter or permission required

State banking authority set operating rules and capital
requirements:
Banks had to post collateral in form of government bonds
Banks could issue their own money (banknotes)
Each banknote dollar had to be collateralized by one dollar
of state or federal bond deposited at the state banking
authority
Banknotes constituted a bank‘s promise to redeem hard
currency (gold or silver coins; called specie) on demand

Banknotes circulated among the population, next to gold coins,
as a form of paper money. But:
They traded at discounts relative to their par value of specie
Discounts, thus, signaled default risk of bank


A wildcat bank is broadly defined as one that prints more currency than it is capable of continuously redeeming in specie. 
A more specific definition, applies the term to free banks whose notes were backed by overvalued securities – 
bonds which were valued at par by the state, but which had a market value below par.

Since the valuation varied from state to state (market value vs. par value), this allowed for the following scheme:
If possible buy discounted bonds for, say, $90 and deposit it with bank authority
Lend out $100 banknotes (=par value of bond) and acquire assets worth $100. 
Sell those assets and close doors of bank 
short term profit of $10 for wildcat bankers and loss for banknote holders (in terms of liquidity).

However, this is seen as unlikely since most bonds traded at par value.

Most bank failures occured when state bond prices fell
The risk of posted collateral (state bonds) made banks
vulnerable to runs
Lesson learned for today:
- Free entry need not be a problem
- Don‘t treat risky collateral as riskless (Italian or Greek Sovereign Bonds?)
Discounts on banknotes discipline banks‘ risk taking, together, reputation and redemption limited wildcat banking.
